\chapter{Ear Fullness}

\section*{Definition}
Ear fullness (aural fullness) is a subjective sensation of pressure, blockage, or congestion in the ear that may be unilateral or bilateral and can occur with or without hearing loss.

\section*{What Happens in Your Body}
The sensation of fullness arises from abnormal pressure in the middle ear (Eustachian tube dysfunction), altered endolymphatic pressure (Meniere disease), stapedius or tensor tympani muscle spasm, or conductive hearing loss that alters the perception of ambient sound pressure levels.

\section*{Causes}
\begin{itemize}
  \item \textbf{Eustachian tube dysfunction:} Most common; mucosal edema from allergy, URI, or sinusitis
  \item \textbf{Middle ear effusion:} Serous otitis media with fluid accumulation
  \item \textbf{Meniere disease:} Endolymphatic hydrops causing episodic aural fullness, vertigo, and fluctuating SNHL
  \item \textbf{Impacted cerumen:} Complete occlusion of the external auditory canal
  \item \textbf{Temporomandibular joint (TMJ) disorders:} Muscle spasm referred to the ear
  \item \textbf{Superior semicircular canal dehiscence:} Third window syndrome with pressure sensitivity
  \item \textbf{Patulous Eustachian tube:} Abnormal patency causing autophony and fullness
\end{itemize}

\section*{When to See a Doctor}
See your doctor if:
\begin{itemize}
  \item Ear fullness persists longer than 2 weeks despite conservative measures
  \item Fluctuating hearing loss, vertigo, or tinnitus accompanies the fullness (suspect Meniere disease)
  \item Unilateral fullness with progressive hearing loss (requires retrocochlear workup)
  \item Fullness associated with otalgia, fever, or otorrhea
  \item Autophony (hearing one's own voice or breath loudly) suggesting patulous tube
\end{itemize}

\section*{Related Symptoms}
\begin{itemize}
  \item Conductive hearing loss (muffled hearing)
  \item Tinnitus (low-frequency roaring in Meniere)
  \item Vertigo or dizziness
  \item Otalgia (ear pain)
  \item Autophony
\end{itemize}
\textbf{Important:} Unilateral aural fullness with unexplained sensorineural hearing loss mandates MRI of the internal auditory canals to exclude a vestibular schwannoma.

\section*{Self-Care / Home Management}
\begin{itemize}
  \item Eustachian tube opening exercises: yawning, swallowing, chewing gum
  \item Oral decongestants or antihistamines for allergic/mucosal congestion
  \item Valsalva maneuver (gentle, holding nose and blowing)
  \item TMJ stretching exercises if jaw tenderness present
  \item Avoiding triggers: caffeine and sodium reduction if Meniere suspected
  \item Warm compress over the ear for muscle tension-related fullness
\end{itemize}

\section*{How Your Doctor Will Evaluate You}
\begin{itemize}
  \item Otoscopy with pneumatic otoscopy to assess TM mobility and middle ear status
  \item Pure-tone audiometry and speech audiometry
  \item Tympanometry to evaluate middle ear pressure and compliance
  \item Tuning fork tests (Rinne and Weber)
  \item MRI of internal auditory canals with gadolinium for asymmetric SNHL
  \item Electrocochleography or glycerol test if Meniere disease suspected
  \item CT temporal bone if superior canal dehiscence in differential
\end{itemize}

\section*{Treatment Options}
\begin{itemize}
  \item \textbf{ET dysfunction:} Topical nasal steroid spray, oral decongestants, and autoinflation; pressure equalization tubes for refractory cases
  \item \textbf{Cerumen impaction:} Cerumenolytic drops (carbamide peroxide) or microsuction removal
  \item \textbf{Meniere disease:} Low-sodium diet (< 2000 mg/day), diuretics, betahistine, intratympanic dexamethasone, or gentamicin ablation
  \item \textbf{Patulous tube:} Hydration, supine positioning, intranasal saline drops; patulous tube occlusion if severe
  \item \textbf{TMJ-related:} Dental splint, NSAIDs, and physical therapy
  \item \textbf{Superior canal dehiscence:} Middle fossa or transmastoid surgical plugging
\end{itemize}

\clearpage
